\documentclass[11pt, a4paper]{article}

\usepackage[T1]{fontenc}
\usepackage[utf8]{inputenc}
\usepackage{tgpagella}
\usepackage{microtype}
\usepackage[spanish, es-noshorthands]{babel}
\addto\captionsspanish{\renewcommand{\tablename}{Tabla}}

\usepackage{amsmath, amssymb, amsfonts}
\usepackage{graphicx}
\usepackage{booktabs}
\usepackage[most]{tcolorbox}
\usepackage[margin=2.5cm]{geometry}

\definecolor{ucbblue}{RGB}{0, 51, 102}

\begin{document}

\begin{center}
    \Large\textbf{\textcolor{ucbblue}{Readiness Check Individual: Cinemática Directa}}\\[0.2cm]
    \normalsize Robótica (IMT-342) - M.Sc. Ing. Bernardo Quiroga Turdera\\[0.1cm]
    \textit{Tiempo objetivo: 10 minutos}[cite: 5, 6]
\end{center}

\vspace{0.4cm}
\textbf{Nombre:} \rule{8cm}{0.15mm} \hfill \textbf{Fecha:} \rule{4cm}{0.15mm}

\vspace{0.6cm}
Considere un manipulador 2R planar con las longitudes de eslabón:
\begin{equation*}
    a_1 = 0.30\text{ m}, \qquad a_2 = 0.20\text{ m}
\end{equation*}

Y su tabla DH estándar:
\begin{center}
    \begin{tabular}{c c c c c}
    \toprule
    $i$ & $\theta_i$ & $d_i$ & $a_i$ & $\alpha_i$ \\ \midrule
    1 & $q_1$ & 0 & 0.30 & 0 \\
    2 & $q_2$ & 0 & 0.20 & 0 \\ \bottomrule
    \end{tabular}
\end{center}

Para las variables articulares $q_1=90^\circ, \ q_2=-90^\circ$ determine[cite: 5]:

\vspace{0.4cm}
\noindent \textbf{1. Matriz de Transformación ${}^0T_1$:}
\vspace{2.5cm}

\noindent \textbf{2. Matriz de Transformación ${}^1T_2$:}
\vspace{2.5cm}

\noindent \textbf{3. Matriz Compuesta ${}^0T_2$:}
\vspace{3.5cm}

\noindent \textbf{4. Posición final $p$:} \rule{6cm}{0.15mm} \\[0.8cm]
\noindent \textbf{5. Orientación final $\phi$ (o matriz $R$):} \rule{6cm}{0.15mm} \\[0.8cm]
\noindent \textbf{6. Verificación (Compruebe geométricamente su resultado):} 

\vfill
\begin{tcolorbox}[colback=gray!10, colframe=gray!50, title=\textbf{Uso Exclusivo del Docente - Scoring[cite: 5]}]
    \small
    $\square$ Usa datos D-H correctos (1) \quad 
    $\square$ ${}^0T_1$ correcto (2) \quad 
    $\square$ ${}^1T_2$ correcto (2) \quad 
    $\square$ Orden de cadena (1) \\
    $\square$ Multiplicación sustancial (2) \quad
    $\square$ Extracción (1) \quad
    $\square$ Verificación (1) \\
    \textbf{Clasificación:} $\square$ Ready (8-10) \quad $\square$ Borderline (5-7) \quad $\square$ Not Yet Ready (0-4)
\end{tcolorbox}

\end{document}
